% Wrapper-free original contribution. Uses the theorem environments and
% macros of bounded_transport.tex. Include after prop:output-crt.
\subsection{The exact incidence between inputs and prescribed prime outputs}

Retain the original prime \(p=12h+1\), \(h\ge1\), and one fixed
positive integer seed
\[
 \varepsilon\in\{0,1\},\qquad A,B,C,D_p>0,\qquad
 4ABCD_p=A+p^{1-\varepsilon}B+pC.
\]
Its original first denominator is \(a_p=ABD_p\); in this
construction retain the full original condition
\(a_p\in A_p=\{3h+1,\ldots,9h\}\).
Keep \(A,B,C,D_p,\varepsilon\) fixed, with
\[
 K_0=4ABC,\quad \Theta=C+(1-\varepsilon)B,\quad
 \Omega=A+\varepsilon B,\quad
 m=K_0/\gcd(K_0,\Theta).
\]
Here \(\Theta\) is the scalar of the preceding proposition.
All other original seed decorations, including any primitive
coordinate conditions, remain attached to this fixed seed.

Retain the input data \(F>0\), \(\lambda_i,\beta_i\in\Z\),
\(m_i>0\) of Proposition~\ref{prop:input-crt}, a positive output
modulus \(\mathcal D\), and integer bounds \(1\le U\le V\).
Run the complete input solvability tests of that proposition.
If they fail, the input set and every incidence set below are empty.
If they pass, denote its computed representative by
\(a_{\rm in}\in\{0,\ldots,M-1\}\), keeping this name distinct
from every first denominator, and write the exact input set as
\begin{align*}
 \mathcal X
 &=\{X\in\Z:U\le X\le V,\ X\equiv0\pmod{6F},\
                  \lambda_iX\equiv\beta_i\pmod{m_i}\ \forall i\}\\
 &=\{a_{\rm in}+Mt:
       \lceil(U-a_{\rm in})/M\rceil\le t\le
       \lfloor(V-a_{\rm in})/M\rfloor\}.
\end{align*}
In particular \(6F\mid M\) and \(6F\mid a_{\rm in}\).
This equality retains every original input congruence; the bounds
are not substituted for its residue conditions.
Every formula below involving \(a_{\rm in},M\) refers to this
computed solvable branch.  When the input tests fail, use
\(\mathcal X=\varnothing\), all incidence fibres empty and all
their counts zero; no residue representative is invented.
Let \(\mathfrak s\) denote this entire fixed labelled datum,
including the seed, original prime and shell, input conditions,
chosen bounds and prescribed output modulus.  Define
\[
 N(X)=X^4-X^2+1,
\]
and the following finite set of candidate output primes:
\[
 \mathcal Q_{\rm out}=
 \{q\text{ prime}:p<q\le N(V),\quad
       q\equiv1\pmod{12},\quad q\equiv1\pmod{\mathcal D},\quad
       q\equiv p\pmod m\}.
\]
The requirement \(q>p\) is an explicit part of this output domain.
For the source specialization
\(F=p!\prod_\nu|E_\nu(p)|\), with every \(E_\nu(p)\ne0\),
the proof below also shows that this requirement is automatic
for every prime divisor of an actual \(N(X)\).

\begin{theorem}[Complete input--output incidence and witness fibres]
\label{thm:input-output-incidence}
The finite incidence set
\[
 \mathcal I_{\mathfrak s}=
 \{(X,q):X\in\mathcal X,\ q\in\mathcal Q_{\rm out},\ q\mid N(X)\}
\]
has the following exact projections, inverses and arithmetic returns.

The projection \(\pi_X:\mathcal I_{\mathfrak s}\to\mathcal X\),
\((X,q)\mapsto X\), has precisely the fibre
\[
 \pi_X^{-1}(X)=
 \left\{(X,q):
 \begin{array}{l}
 q\mid N(X),\quad q\text{ prime},\quad p<q,\\
 q\equiv1\pmod{12},\quad q\equiv1\pmod{\mathcal D},\quad
 q\equiv p\pmod m
 \end{array}\right\}.
\]
Thus its image consists of exactly those bounded inputs whose
displayed filtered prime-divisor list is nonempty.

For the other projection
\(\pi_q:\mathcal I_{\mathfrak s}\to\mathcal Q_{\rm out}\),
\((X,q)\mapsto q\), enumerate, for each candidate \(q\), the full
finite root set
\[
 \mathcal R_q=
 \{\rho\in\{0,\ldots,q-1\}:\rho^4-\rho^2+1\equiv0\pmod q\}.
\]
Put \(g_q=\gcd(M,q)\) and \(L_q=\lcm(M,q)\).
Retain a root exactly when \(g_q\mid\rho-a_{\rm in}\).
For each retained root define
\[
 h_{q,\rho}=
 \operatorname{rem}_{q/g_q}
 \left((M/g_q)^{-1}\frac{\rho-a_{\rm in}}{g_q}\right),\qquad
 x_{q,\rho}=a_{\rm in}+M h_{q,\rho},
\]
where the inverse is modulo \(q/g_q\); when \(q/g_q=1\),
set \(h_{q,\rho}=0\) without taking such an inverse.
The complete integer fibre of the two required input residues is
\(x_{q,\rho}+L_q\Z\).  Its complete bounded branch is
\begin{equation}\label{eq:incidence-root-branch}
 \mathcal T_{q,\rho}=
 \left\{t\in\Z:
    \left\lceil\frac{U-x_{q,\rho}}{L_q}\right\rceil
       \le t\le
    \left\lfloor\frac{V-x_{q,\rho}}{L_q}\right\rfloor\right\}.
\end{equation}
The disjoint union over all these retained roots maps bijectively
onto \(\pi_q^{-1}(q)\) by
\[
 (\rho,t)\longmapsto(x_{q,\rho}+L_qt,q).
\]
Its exact inverse is
\[
 \rho=\operatorname{rem}_q X,\qquad
 t=(X-x_{q,\rho})/L_q.
\]
Each branch has homogeneous congruence kernel \(L_q\Z\).
The incidence projections themselves are maps of finite sets,
with the displayed fibres; no additive kernel is assigned to them.

In particular the full incidence cardinality, including every
bounded lift rather than only one residue representative, is
\begin{align}
 |\mathcal I_{\mathfrak s}|
 &=\sum_{X\in\mathcal X}|\pi_X^{-1}(X)|\notag\\
 &=\sum_{q\in\mathcal Q_{\rm out}}
   \ \sum_{\substack{\rho\in\mathcal R_q\\
                         g_q\mid\rho-a_{\rm in}}}
   \max\left\{0,
       \left\lfloor\frac{V-x_{q,\rho}}{L_q}\right\rfloor
      -\left\lceil\frac{U-x_{q,\rho}}{L_q}\right\rceil+1\right\}.
       \label{eq:incidence-count}
\end{align}
Prime divisors are counted once as labels \(q\), irrespective of
their multiplicity in \(N(X)\).

Every incidence has the explicit original integer witness
\begin{equation}\label{eq:incidence-witness}
 D_q=\frac{\Omega+q\Theta}{4ABC},\qquad
 (a_q,y_q,z_q)=(ABD_q,q^\varepsilon ACD_q,qBCD_q),
 \qquad \frac4q=\frac1{a_q}+\frac1{y_q}+\frac1{z_q}.
\end{equation}
Writing \(h_q=(q-1)/12\), its first denominator lies in the full
original interval
\[
 a_q\in A_q=\{3h_q+1,\ldots,9h_q\}.
\]
Its retained shell coordinates are
\[
 R_q=4a_q-q=\frac{A+q^{1-\varepsilon}B}{C}>0,
 \qquad S_q=qa_q.
\]
The map
\[
 (X,q)\longmapsto
 (\mathfrak s,X,q,D_q,a_q,R_q,S_q,y_q,z_q)
\]
is a bijection onto its explicitly defined graph, with inverse
the projection retaining \(\mathfrak s,X,q\).
The further projection to the labelled output
\((q,D_q,a_q,R_q,S_q,y_q,z_q)\) has exactly the input fibre
described by \eqref{eq:incidence-root-branch} for that \(q\).
Thus no unlabelled inverse silently discards the original input,
prime, seed, or witness ordering.  These maps determine the entire
incidence set and every output witness or prove their exact fibres
empty; they do not assert that either incidence projection is onto.
\end{theorem}
\begin{proof}
First the output search bound is valid for every actual positive
integer input.  Direct expansion gives
\[
 N(1)=1,\qquad
 N(n+1)-N(n)=4n^3+6n^2+2n=2n(n+1)(2n+1)>0
 \quad(n\ge1).
\]
Hence \(1\le X\le V\) implies \(N(X)\le N(V)\).
Every positive prime divisor \(q\) of \(N(X)\) is at most
\(N(X)\), so the stated finite candidate bound loses none.
In fact a retained input is a positive multiple of \(6F\),
so \(X\ge6\) and \(N(X)>1\), although the prescribed prime
filter can still have an empty fibre.

The candidate output list is itself fully computable using the
preceding output CRT.  Put
\(D_*=\lcm(12,\mathcal D)\).
The two requirements with residue one are equivalent to
\(q\equiv1\pmod{D_*}\), because divisibility by both moduli is
divisibility by their least common multiple.
Proposition~\ref{prop:output-crt}, with \(\mathcal D=D_*\),
computes an empty residue set if
\(\gcd(D_*,m)\nmid p-1\), and otherwise computes its complete
residue \(q_0\pmod{L_{\rm out}}\),
\(L_{\rm out}=\lcm(D_*,m)\).
All candidate integers are exactly
\[
 q=q_0+L_{\rm out}v,\qquad
 \left\lceil\frac{p+1-q_0}{L_{\rm out}}\right\rceil
 \le v\le
 \left\lfloor\frac{N(V)-q_0}{L_{\rm out}}\right\rfloor.
\]
Test each for primality to obtain precisely \(\mathcal Q_{\rm out}\).
This test is finite and exact: an integer \(n>1\) is composite
exactly when it has a divisor among
\(2,\ldots,\lfloor\sqrt n\rfloor\).
Indeed a proper factorization \(n=rs\), \(r,s>1\), has at least
one factor no larger than its square root; conversely such a
divisor gives a proper factorization.  The same test, combined
with divisibility by the fixed integer \(N(X)\), computes the
entire filtered prime-divisor fibre of \(\pi_X\).

For the source specialization \(F=p!\prod_\nu|E_\nu(p)|\),
the original input condition \(6F\mid X\) gives
\(N(X)\equiv1\pmod{6F}\).  Every prime at most \(p\)
divides \(p!\), hence \(6F\), and cannot divide \(N(X)\).
Consequently every prime divisor is greater than \(p\).
Likewise no such divisor divides any of the retained nonzero
\(E_\nu(p)\).  This proves the promised exclusion from the
original input divisibility, without adding a prime-generation
or residue-occupancy assertion.  In the general positive \(F\)
case the explicit filter \(q>p\) remains in force.

Fix a candidate output \(q\).  Polynomial evaluation preserves
integer congruences, so
\[
 q\mid N(X)\quad\Longleftrightarrow\quad
 \rho:=\operatorname{rem}_q X\ \text{belongs to }\mathcal R_q.
\]
The left condition \(X\in\mathcal X\) already imposes precisely
\(X\equiv a_{\rm in}\pmod M\) and its interval.
For a chosen root its simultaneous congruence is
\(X\equiv\rho\pmod q\).
The noncoprime CRT compatibility condition is exactly
\(g_q\mid\rho-a_{\rm in}\); an incompatible root gives no
integer lift.  For a compatible root with \(q/g_q>1\),
\(\gcd(M/g_q,q/g_q)=1\), and the displayed inverse and remainder
give a unique integer \(h_{q,\rho}\) in
\(\{0,\ldots,q/g_q-1\}\).
The number \(x_{q,\rho}=a_{\rm in}+Mh_{q,\rho}\) is unchanged
modulo \(M\), and modulo \(q\) the equation for this remainder
gives \(Mh_{q,\rho}\equiv\rho-a_{\rm in}\).
Thus it has both original residues.
When \(q/g_q=1\), compatibility says directly that
\(a_{\rm in}\equiv\rho\pmod q\); the stipulated zero correction
has the same property and \(L_q=M\).

The difference of any two common lifts is divisible by both
\(M\) and \(q\), equivalently by \(L_q\).
Conversely addition of any multiple of \(L_q\) preserves both
residues.  This proves the complete coset and its homogeneous
kernel.  Dividing the unchanged inequalities
\(U\le x_{q,\rho}+L_qt\le V\) by positive \(L_q\) proves
exactly the bounds in \eqref{eq:incidence-root-branch}.
Every such value continues to satisfy \(6F\mid X\) and each
\(\lambda_iX\equiv\beta_i\pmod{m_i}\), because it remains in
the exact input residue class.  None of those conditions has been
dropped after the second CRT merge.

Two different root branches cannot contain the same \(X\):
its residue in \(\{0,\ldots,q-1\}\) is unique.
Inside one branch the positive number \(L_q\) makes its
integer parameter \(t\) unique.  The displayed inverse therefore
recovers both \(\rho\) and \(t\), and every incidence over \(q\)
supplies that inverse.  This proves both directions of the
branch bijection and its disjointness, including empty intervals.
For integers \(b_-,b_+\), the number of integers between them is
\(\max\{0,b_+-b_-+1\}\).  Applying this to each branch and
summing gives the second line of \eqref{eq:incidence-count};
the first counts the same finite labelled pairs by their input
projection.  All searches are finite: candidate primes are
bounded, each root list has \(q\) entries, and every branch
has a finite integer interval.

It remains to prove the complete arithmetic return rather than
identify the output only by a residue class.
The fixed seed gives
\(K_0D_p=\Omega+p\Theta\).
Writing \(g=\gcd(K_0,\Theta)\), \(\Theta=gw\), and
\(K_0=gm\), one has \(\gcd(w,m)=1\).  Hence
\[
 D_q=D_p+\frac{(q-p)w}{m}
\]
is an integer exactly when \(q\equiv p\pmod m\), as proved
in the preceding proposition.  Every retained output satisfies
that congruence, and \(\Omega+q\Theta>0\) proves \(D_q>0\).
Thus all three coordinates in \eqref{eq:incidence-witness} are
positive integers.  For the two retained values of the tag,
\[
 \Omega+q\Theta=A+q^{1-\varepsilon}B+qC.
\]
Putting their three reciprocals over the common positive
denominator \(qABCD_q\) therefore gives
\[
 \frac1{ABD_q}+\frac1{q^\varepsilon ACD_q}
                 +\frac1{qBCD_q}
 =\frac{qC+q^{1-\varepsilon}B+A}{qABCD_q}
 =\frac{4ABCD_q}{qABCD_q}=\frac4q.
\]
No denominator permutation occurs in this calculation.

The original full shell interval also survives on this exact
domain.  Since \(\Theta\ge C\) and \(\Omega>0\),
\[
 \frac{a_q}{q}=\frac{\Theta}{4C}+\frac{\Omega}{4Cq}
                    >\frac14.
\]
As \(q>p\) and \(\Omega>0\), the same exact expression gives
\[
 \frac{a_q}{q}
 \le\frac{a_p}{p}
 \le\frac{3(p-1)}{4p}
 =\frac34-\frac{3}{4p}
 \le\frac34-\frac{3}{4q}.
\]
The second inequality uses precisely the retained original bound
\(a_p\le9h=3(p-1)/4\).
Because \(a_q\) is an integer and \(q=12h_q+1\), the lower
bound \(a_q>q/4\) is equivalent to \(a_q\ge3h_q+1\), and
the upper bound is \(a_q\le3(q-1)/4=9h_q\).
This proves membership in the full original \(A_q\).
Finally \(4a_q=(\Omega+q\Theta)/C\) proves the displayed
formula for \(R_q\).  Its numerator is positive, and
\(R_q=4a_q-q\) proves its integrality; \(S_q=qa_q\) retains
the original supply.  In particular \(R_q\equiv3\pmod4\)
and \(3\le R_q\le2q-3\) follow from that same original
shell interval and \(q\equiv1\pmod{12}\).

The graph construction attaches only coordinates uniquely
computed from the retained \(\mathfrak s,X,q\), so projection
back to those labels is its two-sided inverse and every graph
fibre is a singleton.  On forgetting \(X\) but retaining \(q\)
and the fixed seed, all remaining witness coordinates are fixed;
its full fibre is consequently exactly the previously computed
root-branch union for that \(q\).  This proves the stated
projections, their complete fibres and all integer reconstruction
claims without replacing an input/output incidence by independent
choices of its two coordinates.
\end{proof}

\paragraph{An actual prime divisor excluded by the prescribed class.}
Retain the prime \(p=13\), \(\varepsilon=1\), and
\[
 A=1,\quad B=2,\quad C=1,\quad D_p=2.
\]
The seed identity is \(4\cdot1\cdot2\cdot1\cdot2=16=1+2+13\).
Its original first denominator \(a_p=4\) lies in
\(A_{13}=\{4,\ldots,9\}\), and its ordered witness is
\((4,26,52)\); over denominator \(52\), its reciprocal
numerators are \(13,2,1\), giving \(16/52=4/13\).
Here \(K_0=8\), \(\Theta=1\), \(\Omega=3\), and \(m=8\).
Choose \(F=13!=6\,227\,020\,800\), no further affine input
constraints, and \(\mathcal D=12\).  The exact input residue is
\(a_{\rm in}=0\) modulo \(M=6F=37\,362\,124\,800\).
Set both retained interval endpoints equal to
\[
 U=V=X_0=635\,156\,121\,600=17M.
\]
Thus \(\mathcal X=\{X_0\}\), with its unique original lift
parameter \(17\).  Direct integer division gives
\[
 X_0=73\cdot8\,700\,768\,789+3,
 \qquad N(3)=3^4-3^2+1=73.
\]
Polynomial reduction therefore proves \(73\mid N(X_0)\).
The integer \(73\) is prime: none of \(2,3,5,7\), the primes
at most \(\sqrt{73}<9\), divides it.  It satisfies \(73>13\)
and \(73\equiv1\pmod{12}=\pmod{\mathcal D}\).
However \(73\equiv1\pmod8\), whereas \(13\equiv5\pmod8\).
The exact proposed transport coefficient is
\[
 D_{73}=\frac{\Omega+73\Theta}{K_0}
       =\frac{3+73}{8}=\frac{19}{2},
\]
which is not an integer.  Hence this actual cyclotomic prime
divisor is excluded from \(\mathcal Q_{\rm out}\), and
\((X_0,73)\notin\mathcal I_{\mathfrak s}\).
The unfiltered prime-divisor pair remains explicitly recorded;
the exact incidence tests retain every other divisor that passes
all output constraints.  This fixture proves failure for this
selected divisor and does not assert that the filtered fibre
at \(X_0\) is empty.
